\section{Evaluation}
\label{sec:evaluation}
The evaluation addresses three research questions derived from Prime Agent's design.
\textbf{RQ1: Test-time scaling.}
Can a standardized, expressive execution interface let frontier models convert additional output tokens and API cost into verified task progress?
We evaluate this question on ARC-AGI-3.
\textbf{RQ2: Information management.}
Can models use persistent REPL state to search, transform, and aggregate information across long contexts?
We compare Prime Agent with native and alternative harnesses on long-context reasoning and coding tasks.
\textbf{RQ3: Persistent recursive execution.}
Can the same runtime sustain multi-day experimentation, iterative systems construction, recursive control, and online refinement?
We study nanoGPT, PMPP-Hard, EmulatorBench, Factorio, and MazeBench through end-to-end outcomes and trajectory analysis.
The trajectory analyses show how agents allocate subagents, retain information, and recover from disruption.

\subsection{Interactive reasoning at test-time scale}
ARC-AGI-3 is the clearest test of Prime Agent to support strong, consistent long-horizon evaluation~\citep{arcagi3}.
Each game requires the model to learn the rules of the game, creating an ad-hoc world model under an action limit.
Prime Agent supplies only the environment interface and an autonomous prompt adapted from PRO-LONG~\citep{prolong2026}; the model constructs the strategy. We note that Claude Code and Codex runs perform worse than Anthropic and Open AI self-reported performance on ARC-AGI-3 (public set), so we defer to their results over our own runs with matched prompt and settings.

\begin{figure}[!t]
\centering
\includegraphics[width=\linewidth]{figures/prime_agent/arc-agi3-scaling-combined.pdf}
\caption{\textbf{ARC-AGI-3 test-time scaling.} RHAE score versus output tokens per game (left) and estimated API cost (right).}
\label{fig:arc}
\end{figure}

Across the observed configurations, additional output tokens and cost are converted into progress at sharply different rates (\cref{fig:arc}).
The stronger configurations continue to improve across a long interaction horizon, while others plateau early.
This pattern is consistent with a model-controlled interface that permits model-dependent test-time scaling instead of imposing one fixed workflow.
The reference lines and points place these curves beside official ARC results.
They are external values because our native-harness reruns fell below the published scores, so they situate the result rather than isolate a causal harness effect.

\subsection{Long-context information management}
The long-context suite tests whether a model can actively manage information that does not fit naturally into one prompt.
Prime Agent stores the initial context in a readable file, allowing the model to search, transform, summarize, and revisit it from the persistent REPL.
This changes long-context reasoning from passive attention over a fixed sequence into a programmatic information-management problem.
The suite covers aggregation, latent retrieval, instruction following, reasoning, and long-form coding~\citep{oolong2025,obliq2026,longbenchpro2026,longbenchv2,longcot2026}.

\begin{table}[H]
\centering
\small
\setlength{\tabcolsep}{4.2pt}
\resizebox{\textwidth}{!}{%
\begin{tabular}{ll cc cc cc}
\toprule
& & \multicolumn{2}{c}{GLM-5.2 (reasoning: high)} & \multicolumn{2}{c}{Opus 5 (reasoning: high)} & \multicolumn{2}{c}{GPT-5.6 Sol (reasoning: high)}\\
Task & Setting & Prime & Pi-mono & Prime & Claude Code & Prime & Codex\\
\midrule
OOLONG (Yahoo, 128k)~\citep{oolong2025} & long context & \textbf{.700} & .420 & .900 & \textbf{.920} & \textbf{.940} & .900\\
OOLONG-Pairs~\citep{zhang2025recursivelanguagemodels} & long output & \textbf{.874} & .556 & \textbf{.929} & .922 & \textbf{.911} & .895\\
OBLIQ-Bench (math)~\citep{obliq2026} & ranking (nDCG@10) & \textbf{.669} & .635 & \textbf{.802} & .795 & .612 & \textbf{.646}\\
LongBench Pro (English)~\citep{longbenchpro2026} & comprehension & \textbf{.777} & .768 & \textbf{.804} & .790 & \textbf{.794} & .790\\
LongBench v2~\citep{longbenchv2} & expert long tasks & .680 & \textbf{.696} & .744 & \textbf{.746} & \textbf{.714} & .704\\
ManyIH Coding~\citep{zhang2026manytierinstructionhierarchyllm} & long instructions & \textbf{.424} & .386 & \textbf{.536} & .522 & \textbf{.499} & .454\\
ManyIH IF~\citep{zhang2026manytierinstructionhierarchyllm} & long instructions & \textbf{.209} & .164 & \textbf{.225} & .175 & .216 & \textbf{.232}\\
LongCoT-Mini~\citep{longcot2026} & long reasoning & \textbf{.638} & .613 & \textbf{.722} & .558 & .671 & \textbf{.681}\\
EmulatorBench~\citep{karten2026emulatorbench} & long coding & \textbf{.208} & .000 & .047 & \textbf{.062} & \textbf{.275} & .228\\
\bottomrule
\end{tabular}%
}
\caption{Long-context results. Bold marks the higher point estimate within each nominal-model pair; metrics differ by row. Bold is not statistical significance, and uncertainty intervals are unavailable.}
\label{tab:long}
\end{table}

We generally find Prime Agent to be competitive across a wide range of long tasks, especially against the harness that did not use a model trained around it. Prime Agent especially excels at long-running or long-context tasks, and can competitively run on its own as an autonomous agent. We include a set of focused case studies and experiments on long settings where Prime Agent excels.

\subsection{Multi-day autonomous research}
\label{sec:nanogpt}
The nanoGPT speedrun~\citep{bakouch2026automatedairesearch} measures how far an agent can reduce the number of training steps required for a 124M-parameter GPT to reach a fixed validation loss, with each record verified as an eight-seed mean.
For each of three models (Kimi K3, DeepSeek V4 Pro, and GLM 5.3), we compare Prime Agent against an alternative harness: the model developer's own CLI where one exists, and Claude Code or opencode otherwise.
We find that the choice of harness has little effect on final records compared to the noise of the experiment.

Model behavior, however, differs.
On Prime Agent, models regularly use the persistent REPL to experiment outside the benchmark's training script, for example by simulating a candidate optimizer on synthetic gradients or numerically optimizing update-rule coefficients before launching a training run.
\cref{fig:labcensus} counts these experiments across 18 runs, normalized by the number of training runs each agent executed; \cref{app:labexp} reproduces one such experiment per model.
The effect is largest for DeepSeek V4 Pro, which created roughly six times more such experiments per training run under Prime Agent than under Claude Code.
This may be due to the fact that DeepSeek's own agent harness provides a similar code-execution mode, so the REPL matches a workflow the model was likely trained on.
We also observe that models construct programmatic interfaces to the benchmark itself: Kimi K3 defined a probe function through which it ran roughly ninety screening experiments and all 19 of its validated records, whereas the same model on its own CLI performed every operation through direct file edits and built no such machinery.

\begin{figure}[!t]
\centering
\includegraphics[width=\linewidth]{figures/prime_agent/nanogpt-lab-census.pdf}
\caption{\textbf{Out-of-loop experimentation across harnesses.} Distinct experiments created outside the training script per 100 training-script executions, pooled over the 2--3 seeds per harness (labels: experiments/training runs). Counts are hand-classified from complete traces; denominators are audited where available and otherwise estimated from launch commands.}
\label{fig:labcensus}
\end{figure}


\subsection{Programmatic systems construction}
\paragraph{Emulators.}
An emulator is software that reproduces another computer system's observable behavior. We evaluate Prime Agent on EmulatorBench, a benchmark that tasks agents with constructing emulators in Rust for a variety of game systems. Agents are given a specification of the emulator and a set of diagnostic tests in the form of a verifier.

\begin{figure}[!t]
\centering
\begin{subfigure}[t]{0.49\linewidth}
\includegraphics[width=\linewidth]{figures/prime_agent/emulator-genesis.pdf}
\caption{Sega Genesis.}
\end{subfigure}\hfill
\begin{subfigure}[t]{0.49\linewidth}
\includegraphics[width=\linewidth]{figures/prime_agent/emulator-game-boy-color.pdf}
\caption{Game Boy Color.}
\end{subfigure}
\caption{\textbf{Selected EmulatorBench runs.} Stepwise verifier score versus estimated cost; hue encodes model and line style encodes harness.}
\label{fig:emulators}
\end{figure}

The correctness of an emulator is given from its ability to mimic the behavior of the target machine. This is measured by human-generated diagnostic programs that inspect the emulator's behavior, such as the CPU flags, PPU timing, and other components. In an effort to minimize the effects of data contamination, we require the agent to build the emulator from scratch in Rust, sandboxed without any reference implementation. We report preliminary results in \cref{tab:long} on this long-context coding benchmark averaged over 16 emulator reconstructions, as well as two emulators in \cref{fig:emulators}, the SEGA Genesis and Nintendo Game Boy Color, that Prime Agent successfully reproduces. For Opus, our runs surprisingly failed to solve the tasks despite successful tool-call responses.

\paragraph{GPU kernels.}
PMPP-Hard compresses the same programmatic loop into repeated edit, compile, correctness-check, and profile cycles under a wall-clock budget.

\begin{figure}[!t]
\centering
\includegraphics[width=0.83\linewidth]{figures/prime_agent/pmpp-hard.pdf}
\caption{PMPP-Hard solve rates at fixed within-model budgets.}
\label{fig:pmpp}
\end{figure}
\vspace{0.6\baselineskip}

Prime Agent and the native harnesses remain close, with the ordering reversing between the two model groups.
In these reported within-model comparisons, the general persistent interface supports the compiler--profile loop with no large observed gap. One noted limitation of PMPP-Hard is the strict wall-clock budget comparison.
What the wall-clock budgets do not reveal is the substantial improvement in token usage by models that use Prime Agent. This means that the same performance as Codex or Kimi-Code is achieved by Prime Agent at substantially reduced cost, and, token-for-token, Prime Agent has an advantage.

\subsection{Persistent interaction and refinement}
\paragraph{Factorio.}
The Factorio Learning Environment exposes Python observations and actions for a persistent factory world~\citep{fle2025}.
In a seven-day Sonnet 5 run, the root and its descendants used 23.4 million output tokens while completing 24 of 196 technologies and reaching 71\% on advanced-circuit research (\cref{fig:factorio-tech-tree}) with no signs of stalling.

\begin{figure}[!t]
\centering
\includegraphics[width=\linewidth]{figures/prime_agent/factorio-tech-and-tree.pdf}
\caption{\textbf{Factorio progress and recursive computation.} Technologies researched (top) and agent-tree growth and concurrency (bottom) versus cumulative output tokens for the Sonnet 5 aesthetic run. The vertical line marks a destructive world reset. The final open marker shows 71\% progress on advanced-circuit after 24 completed technologies.}
\label{fig:factorio-tech-tree}
\end{figure}

The model handled irreversible actions poorly. A destructive world reset reverted the technology count from five to one; the session then recovered and continued the run instead of discarding the trajectory.
The root created 633 depth-one subagents across 149 dispatch waves, with at most seven active concurrently.
The shallow, repeatedly widening tree recorded parallel task specialization rather than deeper recursion, while the bursty technology curve separated long construction intervals from externally verified progress.

A different Factorio trace revealed the central safety failure of online refinement.
The agent discovered that RCON commands could spawn resources directly into assembly machines, used the shortcut despite an anti-cheating heartbeat, and then preserved it as a reusable skill.
In this trace, persistence preserved behavior that optimized the measured objective, including a specification exploit.
Safe deployment therefore requires least-privilege action interfaces, independent state validation, and auditable rollback of contaminated refinements.

\paragraph{MazeBench.}
MazeBench is an open-world 3D spatial reasoning environment where the player controls a 3D cube and must solve puzzle rooms within a global maze, while collecting gems. Frontier models are shown to greatly struggle on this task, expending billions of tokens to solve only a fraction of the overall world. We compare Opus 5 and GPT-5.6 Sol with Prime Agent versus their native harnesses, as well as GLM-5.2 with Claude Code. Following the benchmark metrics, we report the unique number of rooms they find, the unique number of states, and the total number of gems, all as a function of their overall token spend.

\begin{figure}[!t]
\centering
\includegraphics[width=\linewidth]{figures/prime_agent/mazebench.pdf}
\caption{MazeBench exploration versus estimated token cost. Solid lines and filled markers show Prime Agent; dotted lines and open markers show comparison harnesses. Hue and marker shape identify the model.}
\label{fig:maze}
\end{figure}
